\QHeading{第1题　在（\hspace{28pt}）里填上适当的数。}
\par\noindent


\PlainItem{$8.9$千克$=$（\textcolor{answerred}{\ensuremath{8}}）千克（\textcolor{answerred}{\ensuremath{900}}）克}
\par

\PlainItem{$3.005$千米$=$（\textcolor{answerred}{\ensuremath{3}}）千米（\textcolor{answerred}{\ensuremath{5}}）米}
\PlainItem{$5$吨$40$千克$=$（\textcolor{answerred}{\ensuremath{5.04}}）吨}
\PlainItem{$1.92$千米$=$（\textcolor{answerred}{\ensuremath{1}}）千米（\textcolor{answerred}{\ensuremath{920}}）米}
\PlainItem{$7$吨$50$千克$=$（\textcolor{answerred}{\ensuremath{7.05}}）吨}
\PlainItem{$3$千克$165$克$=$（\textcolor{answerred}{\ensuremath{3.165}}）千克}
\PlainItem{$5$千米$30$米$=$（\textcolor{answerred}{\ensuremath{5.03}}）千米}
\PlainItem{$8$元$5$角$=$（\textcolor{answerred}{\ensuremath{8.5}}）元}
\par\noindent

\Answer{
\QHeading{计算过程：第1题：复合单位换算}
5吨40千克换算成吨：
\Formula{
\begin{aligned}[t]&\phantom{={}}5+40\div1000\\&=5+0.04\\&=5.04\text{（吨）}\end{aligned}
}
7吨50千克换算成吨：
\Formula{
\begin{aligned}[t]&\phantom{={}}7+50\div1000\\&=7+0.05\\&=7.05\text{（吨）}\end{aligned}
}
3千克165克换算成千克：
\Formula{
\begin{aligned}[t]&\phantom{={}}3+165\div1000\\&=3+0.165\\&=3.165\text{（千克）}\end{aligned}
}
5千米30米换算成千米：
\Formula{
\begin{aligned}[t]&\phantom{={}}5+30\div1000\\&=5+0.03\\&=5.03\text{（千米）}\end{aligned}
}
8元5角换算成元：
\Formula{
\begin{aligned}[t]&\phantom{={}}8+5\div10\\&=8+0.5\\&=8.5\text{（元）}\end{aligned}
}

}

\QHeading{第2题　动手操作。}

\par\noindent\makebox[\linewidth][c]{
\begin{tikzpicture}[x=1bp,y=-1bp,color=black,draw=black,text=black]
\begin{scope}\clip (56.0,286.67) rectangle (557.0,459.33);
\hdhead{95}{总复习}
\hdsection{1}{132}{在（\hspace{28pt}）里填上适当的数。}
\q{95}{184}{375}{12}{$8.9$千克$=$（\hspace{24pt}）千克（\hspace{24pt}）克}
\q{480}{184}{375}{12}{$3.005$千米$=$（\hspace{24pt}）千米（\hspace{24pt}）米}
\q{95}{234}{375}{12}{$5$吨$40$千克$=$（\hspace{32pt}）吨}
\q{480}{234}{375}{12}{$1.92$千米$=$（\hspace{24pt}）千米（\hspace{24pt}）米}
\q{95}{284}{375}{12}{$7$吨$50$千克$=$（\hspace{32pt}）吨}
\q{480}{284}{375}{12}{$3$千克$165$克$=$（\hspace{31pt}）千克}
\q{95}{334}{375}{12}{$5$千米$30$米$=$（\hspace{32pt}）千米}
\q{480}{334}{375}{12}{$8$元$5$角$=$（\hspace{30pt}）元}
\hdsection{2}{397}{动手操作。}
\q{95}{447}{400}{12}{① 把最左边的图形补全，使它成为一}
\q{127}{496}{354}{12}{个轴对称图形。}
\q{95}{545}{390}{12}{② 将平行四边形向右平移5格。}
\q{95}{594}{402}{12}{③ 将三角形绕点$A$顺时针旋转$90^\circ$，}
\q{127}{643}{354}{12}{画出旋转后的图形。}
\draw[hdgrid] (328.000,292.000)--(328.000,404.000);
\draw[hdgrid] (342.000,292.000)--(342.000,404.000);
\draw[hdgrid] (356.000,292.000)--(356.000,404.000);
\draw[hdgrid] (370.000,292.000)--(370.000,404.000);
\draw[hdgrid] (384.000,292.000)--(384.000,404.000);
\draw[hdgrid] (398.000,292.000)--(398.000,404.000);
\draw[hdgrid] (412.000,292.000)--(412.000,404.000);
\draw[hdgrid] (426.000,292.000)--(426.000,404.000);
\draw[hdgrid] (440.000,292.000)--(440.000,404.000);
\draw[hdgrid] (454.000,292.000)--(454.000,404.000);
\draw[hdgrid] (468.000,292.000)--(468.000,404.000);
\draw[hdgrid] (482.000,292.000)--(482.000,404.000);
\draw[hdgrid] (496.000,292.000)--(496.000,404.000);
\draw[hdgrid] (510.000,292.000)--(510.000,404.000);
\draw[hdgrid] (524.000,292.000)--(524.000,404.000);
\draw[hdgrid] (538.000,292.000)--(538.000,404.000);
\draw[hdgrid] (552.000,292.000)--(552.000,404.000);
\draw[hdgrid] (328.000,292.000)--(552.000,292.000);
\draw[hdgrid] (328.000,306.000)--(552.000,306.000);
\draw[hdgrid] (328.000,320.000)--(552.000,320.000);
\draw[hdgrid] (328.000,334.000)--(552.000,334.000);
\draw[hdgrid] (328.000,348.000)--(552.000,348.000);
\draw[hdgrid] (328.000,362.000)--(552.000,362.000);
\draw[hdgrid] (328.000,376.000)--(552.000,376.000);
\draw[hdgrid] (328.000,390.000)--(552.000,390.000);
\draw[hdgrid] (328.000,404.000)--(552.000,404.000);
\draw[dashed,line width=.6pt] (328.000,334.000)--(412.000,334.000);
\draw[line width=.7pt] (342.000,334.000)--(356.000,320.000)--(370.000,334.000)--(384.000,320.000)--(398.000,334.000);
\draw[line width=.7pt] (398.000,334.000)--(384.667,306.667)--(398.000,292.000)--(412.000,306.667)--cycle;
\draw[line width=.7pt] (398.000,390.667)--(426.000,363.333)--(468.000,363.333)--(440.000,390.667)--cycle;
\draw[line width=.7pt] (509.333,306.667)--(537.333,306.667)--(523.333,348.667)--cycle;
\node[inner sep=0pt,align=center,font=\normalsize,] at (503.333,300.000) {$A$};
\hdsection{3}{692}{按要求在下面图形中画一条线段。}
\q{95}{746}{370}{12}{① 分成两个直角梯形。}
\q{478}{746}{370}{12}{② 分成一个平行四边形和一个梯形。}
\draw[line width=.7pt] (132.667,576.667)--(156.000,528.000)--(227.333,528.000)--(204.000,576.667)--cycle;
\draw[line width=.7pt] (369.333,574.000)--(416.667,527.333)--(480.000,527.333)--(500.667,573.333)--cycle;
\hdsection{4}{910}{解决问题。}
\q{95}{959}{745}{12}{王爷爷家兔舍有4层，每层有15只兔子。每只兔子年产毛量达到821}
\q{95}{1009}{746}{12}{克。王爷爷家年产兔毛多少千克？}

\begin{scope}[color=answerred,draw=answerred,text=answerred]% 高清文字版专用填空位置；原页答案另存于 ../overlays。
\ans{140.440}{124.667}{60.000}{11}{$8$}
\ans{206.685}{124.667}{60.000}{11}{$900$}
\ans{159.951}{158.001}{60.000}{11}{$5.04$}
\ans{159.951}{191.334}{60.000}{11}{$7.05$}
\ans{159.951}{224.667}{60.000}{11}{$5.03$}
\ans{408.812}{124.667}{60.000}{11}{$3$}
\ans{480.537}{124.667}{60.000}{11}{$5$}
\ans{402.959}{158.001}{60.000}{11}{$1$}
\ans{469.204}{158.001}{60.000}{11}{$920$}
\ans{419.234}{191.334}{60.000}{11}{$3.165$}
\ans{400.550}{224.667}{60.000}{11}{$8.5$}
\begin{scope}[scale=0.6666667]\draw[answerred,line width=1pt] (513,501)--(534,522)--(555,501)--(576,522)--(597,501);\draw[answerred,line width=1pt] (597,501)--(577,542)--(598,564)--(618,542)--cycle;\draw[answerred,line width=1pt] (702,586)--(744,545)--(806,545)--(764,586)--cycle;\draw[answerred,line width=1pt] (764,460)--(764,502)--(702,481)--cycle;\node[answerred,font=\normalsize] at (514,534) {①};\node[answerred,font=\normalsize] at (750,596) {②};\node[answerred,font=\normalsize] at (730,459) {③};\draw[answerred,line width=1pt] (270,792)--(270,865);\draw (270,800)--(278,800)--(278,792);\draw (270,857)--(278,857)--(278,865);\draw[answerred,line width=1pt] (673,791)--(705,860);\end{scope}
\ans{86.667}{697.000}{466.667}{10.5}{{\setlength{\jot}{-2pt}$\begin{aligned}[t]&\phantom{={}}821\times(15\times4)\div1000\\&=821\times60\div1000\\&=49260\div1000\\&=49.26\text{（千克）}\end{aligned}$}\\答：全年产兔毛49.26千克。}
\ans{84.000}{787.000}{460.000}{9}{完整计算过程见紧接本页的计算续页。}
\end{scope}
\end{scope}\pgfresetboundingbox\path[use as bounding box] (56.0,286.67) rectangle (557.0,459.33);
\end{tikzpicture}}\par



\par



\par\noindent

\QHeading{第3题　按要求在下面图形中画一条线段。}

\par\noindent\makebox[\linewidth][c]{
\begin{tikzpicture}[x=1bp,y=-1bp,color=black,draw=black,text=black]
\begin{scope}\clip (56.0,492.0) rectangle (508.0,604.67);
\hdhead{95}{总复习}
\hdsection{1}{132}{在（\hspace{28pt}）里填上适当的数。}
\q{95}{184}{375}{12}{$8.9$千克$=$（\hspace{24pt}）千克（\hspace{24pt}）克}
\q{480}{184}{375}{12}{$3.005$千米$=$（\hspace{24pt}）千米（\hspace{24pt}）米}
\q{95}{234}{375}{12}{$5$吨$40$千克$=$（\hspace{32pt}）吨}
\q{480}{234}{375}{12}{$1.92$千米$=$（\hspace{24pt}）千米（\hspace{24pt}）米}
\q{95}{284}{375}{12}{$7$吨$50$千克$=$（\hspace{32pt}）吨}
\q{480}{284}{375}{12}{$3$千克$165$克$=$（\hspace{31pt}）千克}
\q{95}{334}{375}{12}{$5$千米$30$米$=$（\hspace{32pt}）千米}
\q{480}{334}{375}{12}{$8$元$5$角$=$（\hspace{30pt}）元}
\hdsection{2}{397}{动手操作。}
\q{95}{447}{400}{12}{① 把最左边的图形补全，使它成为一}
\q{127}{496}{354}{12}{个轴对称图形。}
\q{95}{545}{390}{12}{② 将平行四边形向右平移5格。}
\q{95}{594}{402}{12}{③ 将三角形绕点$A$顺时针旋转$90^\circ$，}
\q{127}{643}{354}{12}{画出旋转后的图形。}
\draw[hdgrid] (328.000,292.000)--(328.000,404.000);
\draw[hdgrid] (342.000,292.000)--(342.000,404.000);
\draw[hdgrid] (356.000,292.000)--(356.000,404.000);
\draw[hdgrid] (370.000,292.000)--(370.000,404.000);
\draw[hdgrid] (384.000,292.000)--(384.000,404.000);
\draw[hdgrid] (398.000,292.000)--(398.000,404.000);
\draw[hdgrid] (412.000,292.000)--(412.000,404.000);
\draw[hdgrid] (426.000,292.000)--(426.000,404.000);
\draw[hdgrid] (440.000,292.000)--(440.000,404.000);
\draw[hdgrid] (454.000,292.000)--(454.000,404.000);
\draw[hdgrid] (468.000,292.000)--(468.000,404.000);
\draw[hdgrid] (482.000,292.000)--(482.000,404.000);
\draw[hdgrid] (496.000,292.000)--(496.000,404.000);
\draw[hdgrid] (510.000,292.000)--(510.000,404.000);
\draw[hdgrid] (524.000,292.000)--(524.000,404.000);
\draw[hdgrid] (538.000,292.000)--(538.000,404.000);
\draw[hdgrid] (552.000,292.000)--(552.000,404.000);
\draw[hdgrid] (328.000,292.000)--(552.000,292.000);
\draw[hdgrid] (328.000,306.000)--(552.000,306.000);
\draw[hdgrid] (328.000,320.000)--(552.000,320.000);
\draw[hdgrid] (328.000,334.000)--(552.000,334.000);
\draw[hdgrid] (328.000,348.000)--(552.000,348.000);
\draw[hdgrid] (328.000,362.000)--(552.000,362.000);
\draw[hdgrid] (328.000,376.000)--(552.000,376.000);
\draw[hdgrid] (328.000,390.000)--(552.000,390.000);
\draw[hdgrid] (328.000,404.000)--(552.000,404.000);
\draw[dashed,line width=.6pt] (328.000,334.000)--(412.000,334.000);
\draw[line width=.7pt] (342.000,334.000)--(356.000,320.000)--(370.000,334.000)--(384.000,320.000)--(398.000,334.000);
\draw[line width=.7pt] (398.000,334.000)--(384.667,306.667)--(398.000,292.000)--(412.000,306.667)--cycle;
\draw[line width=.7pt] (398.000,390.667)--(426.000,363.333)--(468.000,363.333)--(440.000,390.667)--cycle;
\draw[line width=.7pt] (509.333,306.667)--(537.333,306.667)--(523.333,348.667)--cycle;
\node[inner sep=0pt,align=center,font=\normalsize,] at (503.333,300.000) {$A$};
\hdsection{3}{692}{按要求在下面图形中画一条线段。}
\q{95}{746}{370}{12}{① 分成两个直角梯形。}
\q{478}{746}{370}{12}{② 分成一个平行四边形和一个梯形。}
\draw[line width=.7pt] (132.667,576.667)--(156.000,528.000)--(227.333,528.000)--(204.000,576.667)--cycle;
\draw[line width=.7pt] (369.333,574.000)--(416.667,527.333)--(480.000,527.333)--(500.667,573.333)--cycle;
\hdsection{4}{910}{解决问题。}
\q{95}{959}{745}{12}{王爷爷家兔舍有4层，每层有15只兔子。每只兔子年产毛量达到821}
\q{95}{1009}{746}{12}{克。王爷爷家年产兔毛多少千克？}

\begin{scope}[color=answerred,draw=answerred,text=answerred]% 高清文字版专用填空位置；原页答案另存于 ../overlays。
\ans{140.440}{124.667}{60.000}{11}{$8$}
\ans{206.685}{124.667}{60.000}{11}{$900$}
\ans{159.951}{158.001}{60.000}{11}{$5.04$}
\ans{159.951}{191.334}{60.000}{11}{$7.05$}
\ans{159.951}{224.667}{60.000}{11}{$5.03$}
\ans{408.812}{124.667}{60.000}{11}{$3$}
\ans{480.537}{124.667}{60.000}{11}{$5$}
\ans{402.959}{158.001}{60.000}{11}{$1$}
\ans{469.204}{158.001}{60.000}{11}{$920$}
\ans{419.234}{191.334}{60.000}{11}{$3.165$}
\ans{400.550}{224.667}{60.000}{11}{$8.5$}
\begin{scope}[scale=0.6666667]\draw[answerred,line width=1pt] (513,501)--(534,522)--(555,501)--(576,522)--(597,501);\draw[answerred,line width=1pt] (597,501)--(577,542)--(598,564)--(618,542)--cycle;\draw[answerred,line width=1pt] (702,586)--(744,545)--(806,545)--(764,586)--cycle;\draw[answerred,line width=1pt] (764,460)--(764,502)--(702,481)--cycle;\node[answerred,font=\normalsize] at (514,534) {①};\node[answerred,font=\normalsize] at (750,596) {②};\node[answerred,font=\normalsize] at (730,459) {③};\draw[answerred,line width=1pt] (270,792)--(270,865);\draw (270,800)--(278,800)--(278,792);\draw (270,857)--(278,857)--(278,865);\draw[answerred,line width=1pt] (673,791)--(705,860);\end{scope}
\ans{86.667}{697.000}{466.667}{10.5}{{\setlength{\jot}{-2pt}$\begin{aligned}[t]&\phantom{={}}821\times(15\times4)\div1000\\&=821\times60\div1000\\&=49260\div1000\\&=49.26\text{（千克）}\end{aligned}$}\\答：全年产兔毛49.26千克。}
\ans{84.000}{787.000}{460.000}{9}{完整计算过程见紧接本页的计算续页。}
\end{scope}
\end{scope}\pgfresetboundingbox\path[use as bounding box] (56.0,492.0) rectangle (508.0,604.67);
\end{tikzpicture}}\par



\par



\par\noindent

\QHeading{第3题①　平行四边形的两种分法}

\Answer{分别画两条对角线中的一条，每次都能分成两个三角形。下面是两种独立答案，不是在同一个图中把两条线都画出来。}
\par\smallskip\noindent\begin{adjustbox}{max width=\linewidth,center}
\begin{tikzpicture}[x=.65cm,y=.65cm]
\draw[black,thick] (0,0)--(4,0)--(5,2)--(1,2)--cycle;\draw[answerred,thick] (0,0)--(5,2);
\begin{scope}[xshift=4.5cm]\draw[black,thick] (0,0)--(4,0)--(5,2)--(1,2)--cycle;\draw[answerred,thick] (1,2)--(4,0);\end{scope}
\end{tikzpicture}
\end{adjustbox}\par\smallskip

\QHeading{第3题②　梯形分割的另一种画法}

\Answer{除原图给出的分割线外，也可在较长底边上另选一个合适的点，过这个点作平行于某一条腰的线段，与另一条底边相交，分出平行四边形和梯形。交点应落在另一底边内部，不能落在端点而把剩下的一块分成三角形。}

\QHeading{第4题　解决问题。}

\Problem{王爷爷家兔舍有4层，每层有15只兔子。每只兔子年产毛量达到821克。王爷爷家年产兔毛多少千克？}
\par


\MethodTitle{方法一}
\Answer{{\Formula{\begin{aligned}[t]&\phantom{={}}821\times(15\times4)\div1000\\&=821\times60\div1000\\&=49260\div1000\\&=49.26\text{（千克）}\end{aligned}}}\par 答：全年产兔毛49.26千克。}


\MethodTitle{方法二：先求一层兔子的年产毛量}
\Answer{先算一层15只兔子，再乘4层，最后把克换成千克。}
\Answer{\Formula{\begin{aligned}[t]&\phantom{={}}821\times 15\times 4\div 1000\\
&=12315\times 4\div 1000\\
&=49260\div 1000\\
&=49.26\text{（千克）}\end{aligned}}}
\Answer{答：年产兔毛49.26千克。}


